% !TeX program = xelatex
\documentclass{article}
\usepackage{kotex}
\usepackage[a4paper, left=3cm, right=3cm, top=2.5cm, bottom=2.5cm]{geometry}
\usepackage{amsmath, amssymb}
\usepackage{tcolorbox}
\usepackage{caption}
\usepackage{setspace}
\usepackage{xcolor}
\usepackage{titlesec}
\usepackage{tikz}
\usetikzlibrary{arrows.meta}

\setstretch{1.3}
\renewcommand{\figurename}{그림}
\titleformat{name=\section,numberless}
  {\normalfont\large\bfseries\color{blue}}{}
  {0pt}{}

% 유기적 실루엣 (높이≈1.2)
\newcommand{\drawbird}{%
  \draw[fill=black!72, draw=none]
    plot[smooth cycle, tension=0.70] coordinates {%
      (0.00,1.20)(-0.20,1.05)(-0.38,0.70)(-0.30,0.30)%
      (-0.05,0.00)(0.20,0.05)(0.28,0.40)(0.18,0.85)%
    };%
}
% 동일 실루엣 ×1.45 확대
\newcommand{\drawbigbird}{%
  \draw[fill=black!72, draw=none]
    plot[smooth cycle, tension=0.70] coordinates {%
      (0.00,1.74)(-0.29,1.52)(-0.55,1.02)(-0.44,0.44)%
      (-0.07,0.00)(0.29,0.07)(0.41,0.58)(0.26,1.23)%
    };%
}

\title{1.2 등거리 변환과 합동}
\date{}
\author{}

\begin{document}
\maketitle

\section*{1.2 등거리 변환과 합동}

\begin{tcolorbox}[
  colback=teal!4, colframe=teal!65!black,
  title=정의 \& 핵심 규칙,
  fonttitle=\bfseries, boxrule=0.5mm, arc=3mm]

\textbf{\textcolor{teal!65!black}{정의 1.2.1 (등거리 변환).}}\quad
함수 $f : \mathbb{E}^n \to \mathbb{E}^n$이 임의의 두 점 $P, Q$에 대해
\[
  f(P)\,f(Q) \;=\; PQ
\]
를 만족하면 $f$를 \textbf{등거리 변환(isometry)}이라 한다.
(유클리드 등거리 변환, 유클리드 운동, 유클리드 변환이라고도 부른다.)

\noindent\textcolor{teal!50!black}{\rule{\linewidth}{0.4pt}}

\textbf{\textcolor{teal!65!black}{핵심 규칙 (보존량).}}\quad
등거리 변환은 거리를 보존하므로, 다음을 모두 보존한다.

\smallskip
\begin{center}
\renewcommand{\arraystretch}{1.25}
\begin{tabular}{|l|l|}
\hline
보존되는 양 & 이유 \\
\hline
길이(거리) & 정의에 의해 직접 보존 \\
각 & 길이로부터 코사인 법칙으로 유도 \\
넓이·부피 & 길이와 각으로부터 유도 \\
\hline
\end{tabular}
\end{center}

\smallskip
$\Rightarrow$ 등거리 변환은 모든 기하학적 도형의 \textbf{크기와 모양}을 보존한다.

\end{tcolorbox}

\medskip
\textbf{보기 1.2.1.}\quad
선분 $\overline{PQ}$는 $P$와 $Q$를 잇는 가장 짧은 곡선이다.
$f$가 등거리 변환이면 거리를 보존하므로 $f(\overline{PQ})$도 $f(P)$와 $f(Q)$를 잇는
가장 짧은 곡선, 즉 선분이다. 따라서
\[
  f\!\left(\overline{PQ}\right) = \overline{f(P)\,f(Q)}
\]

\begin{figure}[h]
\centering
\begin{tikzpicture}[>=Stealth]
  % === 왼쪽: 등거리 변환 f ===
  \drawbird
  \node[above, inner sep=1pt] at ( 0.00, 1.27) {\footnotesize $P$};
  \node[below, inner sep=1pt] at (-0.05, 0.00) {\footnotesize $Q$};

  \draw[->, thick] (0.60,0.60) -- (1.90,0.60) node[midway,above]{$f$};

  \begin{scope}[xshift=2.55cm]
    \drawbird
    \node[above, inner sep=1pt] at ( 0.00, 1.27) {\footnotesize $f(P)$};
    \node[below, inner sep=1pt] at (-0.05, 0.00) {\footnotesize $f(Q)$};
  \end{scope}

  \node[below] at (1.27,-0.38) {등거리 변환};

  % === 구분선 ===
  \draw[dashed, gray!60, line width=0.5pt] (3.58,-0.52) -- (3.58,1.68);

  % === 오른쪽: 등거리 변환 아님 g ===
  \begin{scope}[xshift=4.18cm]
    \drawbird
    \node[above, inner sep=1pt] at ( 0.00, 1.27) {\footnotesize $P$};
    \node[below, inner sep=1pt] at (-0.05, 0.00) {\footnotesize $Q$};

    \draw[->, thick] (0.60,0.60) -- (1.90,0.60) node[midway,above]{$g$};

    \begin{scope}[xshift=2.55cm]
      \drawbigbird
      \node[above, inner sep=1pt] at ( 0.00, 1.82) {\footnotesize $g(P)$};
      \node[below, inner sep=1pt] at (-0.07, 0.00) {\footnotesize $g(Q)$};
    \end{scope}

    \node[below] at (1.27,-0.38) {등거리 변환이 아니다};
  \end{scope}
\end{tikzpicture}
\caption{유클리드 공간에서의 함수 ($f$는 등거리 변환, $g$는 아님)}
\label{fig:1.1}
\end{figure}

\begin{tcolorbox}[
  colback=teal!4, colframe=teal!65!black,
  title=정의 1.2.2\quad 합동,
  fonttitle=\bfseries, boxrule=0.5mm, arc=3mm]

두 부분집합 $A, B \subset \mathbb{E}^n$에 대해
$f(A) = B$를 만족하는 등거리 변환 $f$가 존재하면,
$A$와 $B$는 \textbf{합동(congruent)}이라 하고
\[
  A \;\cong\; B
\]
로 나타낸다.

\end{tcolorbox}

\medskip
\subsection*{연습문제 1.2}

\begin{enumerate}
  \item \textbf{(문제 1.2.1)}
    \begin{enumerate}
      \item[(a)] 등거리 변환은 단사(일대일) 함수임을 보여라.\\
        (즉, $f$가 등거리 변환이고 $P \neq Q$이면 $f(P) \neq f(Q)$임을 보여라.)
      \item[(b)] $f$가 등거리 변환이고 역함수 $f^{-1}$이 존재하면, $f^{-1}$도 등거리 변환임을 보여라.
    \end{enumerate}

  \item \textbf{(문제 1.2.2)}\quad
    $C \subset \mathbb{E}^2$를 중심 $P$, 반지름 $r$인 원이라 하자.
    등거리 변환 $f : \mathbb{E}^2 \to \mathbb{E}^2$에 대해
    $f(C)$가 중심 $f(P)$, 반지름 $r$인 원임을 증명하여라.\\
    \textit{(도움말: 중심 $f(P)$, 반지름 $r$인 원을 $C'$이라 하고
    $f(C) \subset C'$와 $f^{-1}(C') \subset C$가 성립함을 보여라.)}
\end{enumerate}

\end{document}
